\documentclass[addpoints]{exam}
% \documentclass[addpoints,12pt]{exam}

\usepackage[slovene]{babel}
\usepackage{amsfonts}
\usepackage{amssymb}
\usepackage{amsmath}
\usepackage{float}
\usepackage{graphicx}
\usepackage{enumitem}
\usepackage{color}
\usepackage{eurosym}


%Komentar naslednje vrstice glede na verzijo testa -- rešitve ali prostor za odgovore:
% \printanswers

% \linespread{2}


\pointsinrightmargin
\marginbonuspointname{}


\renewcommand{\solutiontitle}{\noindent\textbf{Rešitve in točkovanje:}\par\noindent}

\colorgrids
\definecolor{GridColor}{gray}{0.7}
\setlength{\gridsize}{7mm}

\extrawidth{5mm}
\setlength{\rightpointsmargin}{1.5cm}

\begin{document}

\runningfooter{}{}{}

\begin{center}
    \fbox{\fbox{\parbox{15cm}{\centering
    Popravljanje in pridobivanje ocen -- konferenca \\ 1. letnik -- test C \\ 17. junij 2025 \\ (Osnove logike in teorije množic; Naravna in cela števila; Potence in izrazi; Deljivost)}}}
\end{center}

\vspace{0.1in}
\makebox[\textwidth]{Ime in priimek, oddelek:\enspace\hrulefill}


\begin{center}
    \hqword{Naloga}
    \hpword{Možne točke}
    \chbpword{Dodatne točke}
    \hsword{Dosežene točke}
    \htword{Skupaj}  
    % \combinedpointtable[h][questions]
    \gradetable[h][questions]

\end{center}

\vspace{0.1in}


\begin{questions}

    % 1. vprašanje
    
    \question[4]
    Določite pravilnost izjave $ A\Leftrightarrow B\Rightarrow (\lnot A\land B)$ glede na izbor izjav $A$ in $B$. Upoštevajte vrstni red operacij.

    \begin{solution}[12cm]
        \begin{table}[H]
            \begin{tabular}{ccccccc}
             $A$ & $B$ & $\lnot A$ & $A\Rightarrow B$ & $\lnot A\land B$ & $\lnot(\lnot A\land B)$ & $(A\Rightarrow  B)\Leftrightarrow \lnot(A\land \lnot B)$\\ \hline
             P & P & N & P & N & P & P \\ 
             P & N & N & N & N & P & N \\
             N & P & P & P & P & N & N \\
             N & N & P & P & N & P & P
            \end{tabular}
        \end{table}
        Za vsako pravilno izpolnjeno vrstico, ki določa pravilnost, po $1$ točka.

    \end{solution}




    % 2. vprašanje
    \question[2]
    Določite resničnostno vrednost izjav:
    \begin{parts}
        % \part 
        % Število $5$ je prvo liho praštevilo natanko tedaj, ko število $9$ ni praštevilo.\\
        \part 
        Če število $3$ ni praštevilo, potem je praštevil števil števno končno.\\
        \part
        Ni res, da je ostanek pri deljenju s $5$ lahko samo $1$, $2$ ali $0$, in število $7$ ni večje ali enako številu $8$.\\
    \end{parts}

    \begin{solution}
        (a): P (prva izjava je nepravilna, druga je nepravilna) \\
        (b): N (negirana prva izjava je nepravilna /prva izjava je pravilna/, druga izjava je nepravilna) \\ \\
        Za vsako pravilno določeno logično vrednost po $1$ točka.
    \end{solution}



    \newpage

    % 3. vprašanje
    \question
    Naj bo $\mathcal{U}=\mathbb{N}_{12}$ ter $\mathcal{A}=\left\{x;\ x~ je~ delitelj~ 9\right\}$ in $\mathcal{B}=\left\{x;\ x=2n+1\land n\in\mathbb{N}\right\}$.
    \begin{parts}
        \part[3]
        Zapišite dane množice z elementi, določite njihovo moč. 
        \part[6]
        Določite naslednje množice: $\mathcal{C}=\mathcal{A}\cap\mathcal{B}$, $\mathcal{D}=\mathcal{A}\setminus\mathcal{B}$, $\mathcal{B}\times\mathcal{D}$, $\mathcal{PC}$ in $(\mathcal{A}\cap\mathcal{B})^\complement$. Kartezični produkt predstavite grafično.
        % \part[1]
        % Zapišite najmanjšo možno množico $\mathcal{E}$, da bosta množici $\mathcal{A}$ in $\mathcal{B}$ njeni podmnožici.
        \part[1]
        Zapišite množico $\mathcal{F}$, ki je disjunktna z množico $\mathcal{A}$ in velja $\mathcal{F}\subseteq\mathcal{B}$ ter $m(\mathcal{F})=2$.
    \end{parts}


    \begin{solution}[12cm]
        (a): $\mathcal{A}=\{1,2,3,6\}$, $\mathcal{B}=\{2,4,6,8,10\}$ in $\mathcal{U}=\{1,2,3,4,5,6,7,8,9,10,11\}$; $m(\mathcal{A})=4$, $m(\mathcal{B})=5$, $m(\mathcal{U})=11$  \\
        (b): $\mathcal{C}=\mathcal{A}\cap\mathcal{B}=\{2,6\}$, $\mathcal{A}\cup\mathcal{B}=\{1,2,3,4,6,8,10\}$, $\mathcal{D}=\mathcal{A}\setminus\mathcal{B}=\{1,3\}$, $\mathcal{PC}=\{\emptyset,\{2\},\{6\},\mathcal{C}\}$, $\mathcal{B}\times\mathcal{D}=\{(2,1),(2,3),(4,1),(4,3),(6,1),(6,3),(8,1),(8,3),(10,1),(10,3)\}$ in $(\mathcal{A}\cup\mathcal{B})^\complement=\{5,7,9,11\}$ + grafična upodobitev\\
        (c): $\mathcal{E}=\{1,2,3,4,6,8,10\}$ \\
        Pri (a) za zapis množic $1$ točka; za zapis moči množic $1$ točka. Pri (b) za zapis vsake množice po $1$ točka, za grafično upodobitev $1$ točka. Pri (c) za zapis množice $1$ točka.
    \end{solution}






      \newpage
    % 5. vprašanje
    \question
    Dijaki in dijakinje 3. F razreda so ugotavljali, katero znamko računalnika imajo.
    Računalnik znamke \textit{IBM} ima doma $13$ dijakov, znamke \textit{Dell} $11$ dijakov, znamke \textit{Apple} pa tudi $11$ dijakov.
    Vse tri znamke računalnika imata doma $2$ dijaka. Računalnik znamke \textit{Apple} in \textit{Dell} $4$ dijaki, \textit{Apple} in \textit{IBM} $3$ dijaki, \textit{IBM} in \textit{Dell} pa $5$ dijakov.
    Trije dijaki imajo doma računalnik neke druge znamke.
    \begin{parts}
        \part[1]
        Narišite diagram.
        \part[2]
        Koliko je vseh dijakov v razredu?
        \part[2]
        Koliko je takih dijakov, ki imajo doma natanko en računalnik?
    \end{parts}

    \begin{solution}[10cm]
        (a): Euler-Vennov diagram z vpisanimi številkami \\
        (b): $110-5-49-27=29$ gospodinjstev\\
        (c): $5+46+27+49=127$ oziroma $5+46+27=78$ gospodinjstev \\ \\
        Pri (a) za pravilno narisan in izpolnjen diagram $1$ točka. Pri (b) in (c) za pravilen izračun $1$ točka, za zapis odgovora $1$ točka.
    \end{solution}




    %  \newpage
        

    % 6. vprašanje
    % \question[3]
    % Na šoli se izvaja pet različnih krožkov. Pri naravoslovne so tri skupine po $22$ učencev, pri astronmskem je šest skupin z $14$ učenci. Pri odbojki je oblikovanih $8$ skupin s po $7$ učenci. Šah obiskuje $32$ učencev v $5$ skupinah. Popotniški krožek pa ima dve skupini s po $146$ učenci.
    % Med vsemi učenci je 320 fantov. Koliko deklet je na šoli vključenih v krožke? (Zapišite samo en račun.)

    % \begin{solution}[8cm]
    %     $3\cdot 22+6\cdot 14+8\cdot 7+32+2\cdot 146-320=210$ deklet \\ \\
    %     Za nastavek enega računa $1$ točka, za pravilen izračun $1$ točka, za zapis odgovora $1$ točka.
    % \end{solution}

    
    % %dodatno vprašanje
    % \bonusquestion[3]
    % \textit{Dodatna naloga:} Rešite enačbo $(X\setminus\{b\})\times(X \cap \{a,b\})=\left\{(a,a),(c,a),(d,a),(e,a)\right\}$.

    % \begin{solution}
    %     $X=\left\{a,c,d,e\right\}$ \\
    %     Za v celoti pravilen zapis rešitve $3$ točke.
    % \end{solution}










    % 1. vprašanje
    
    \question
    Izračunajte.
\begin{parts}
    % \part[3]
    %     $b^2 \left(a^2+b^2\right) -b^2-\left(-b\right)^2a^2$

    %     \begin{solution}[5cm]
    %         $$b^2 \left(a^2+b^2\right) -b^2-\left(-b\right)^2a^2=b^2a^2+b^4-b^2-b^2a^2=b^4-b^2=b^2(b^2-1)$$
    %         Za pravilno potenciranje negativnega predznaka $1$ točka, pravilno množenje $1$ točka, končen rezultat $1$ točka.
    %     \end{solution}
        
        
    \part[4]
    $\left(-1\right)^{2009}-\left(\left(-4\right)^3-5^3+\left(3^2+4\right)^2\right)^2$

    \begin{solution}[5cm]
        \begin{align*}
            \left(-1\right)^{2025}+\left(-\left(-3\right)^4+5^3+\left(6-3^2\right)^3\right)^2&=-1+(-81+125+(6-9)^3)^2= \\
                                                                                            &=-1+(44+(-3)^3)^2= \\
                                                                                            &=-1+(44+(-27))^2= \\
                                                                                            &=-1+289= \\
                                                                                            &=288
        \end{align*}
        Za pravilno uporabo pravila potenciranja $-1$ $1$ točka, pravilno potenciranje z negativnim predznakom $1$ točka, pravilno množenje $1$ točka, končen rezultat $1$ točka.
    \end{solution}


    
    \part[3]
        $b^2 \left(a^2+b\right) -b^3-\left(-b^2\right)a^2$

    \begin{solution}[2cm]
        $$4\cdot x^{n+3}-5x^n x^3+2x\cdot x^{n+2}=4x^{n+3}-5x^{n+3}+2x^{n+3}=x^{n+3}(4-5+2)=x^{n+3}$$
        Za pravilen zapis potenc števila $x$ $1$ točka, pravilno izpostavljanje skupnega faktorja $1$ točka, končen rezultat $1$ točka.
    \end{solution}





        
    \end{parts}




     \newpage

    % 2. vprašanje
    % \question[3]
    % Izračunajte na čim krajši način: $444444443^2-444444442^2$.   

    % \begin{solution}[5cm]
    %     $$444444443^2-444444442^2=(444444443-444444442)(444444443+444444442)=1\cdot 444444445=444444445$$
    %     Za pravilno uporabo pravila razlike kvadratov $1$ točka, za pravilno množenje $1$ točka, končen rezultat $1$ točka.
    % \end{solution}




    % 3. vprašanje
        \question[3]
        Izračunajte na čim krajši način: $(2312+50)^2-(2312-50)^2$.   
    
        \begin{solution}[8.5cm]
            \begin{align*}
                &(a-2b)(a+2b)-(a-b)^2+(a+b)^2+8b^2= \\
                =&a^2-4b^2-a^2+2ab-b^2+a^2+2ab+b^2+8b^2= \\
                =&a^2+4ab+4b^2= \\
                =&(a+2b)^2
            \end{align*}
            Za pravilno kvadriranje binoma $1$ točka, pravilno množenje in seštevanje izrazov $1$ točka, pravilno izpostavljanje $1$ točka, končen rezultat $1$ točka.
        \end{solution}
    

    % \newpage
    % 4. vprašanje
    \question
    Faktorizirajte izraze, kolikor je mogoče.
    \begin{parts}

        \part[2]
        $v^2u^2-256u^2$

        \begin{solution}[4cm]
            $$9x^2-x^3+9x-81=x^2(9-x)-9(-x+9)=(x^2-9)(9-x)=(x-3)(x+3)(9-x)$$
            Za uporabo pravila faktorizacije štiričlenika 2+2 $1$ točka, pravilna faktorizacija štiričlenika $1$ točka, končen rezultat $1$ točka.
        \end{solution}

        % \part[2]
        % $a^2c-169c$

        % \begin{solution}[4cm]
        %     $$a^2c-169c=c(a^2-169)=c(a-13)(a+13)$$
        %     Za pravilno izpostavljanje $1$ točka, pravilno uporabljena formula razlike kvadratov $1$ točka.
        % \end{solution}

        \part[2]
        $x^2-11xz^2+28z^4$

        \begin{solution}[4cm]
            $$g^2-13gh+40h^2=(g-5h)(g-8h)$$
            Za uporabo Vietovega pravila $1$ točka, pravilen rezultat $1$ točka.
        \end{solution}


        % \part[3]
        % $x^2+4xy+4y^2-81$

        % \begin{solution}[4cm]
        %     $$x^2+4xy+4y^2-81=(x+2y)^2-9^2=(x+2y-9)(x+2y+9)$$
        %     Za uporabo pravila faktorizacije štiričlenika 3+1 $1$ točka, pravilna uporaba formule $1$ točka, pravilno razstavljanje razlike kvadratov $1$ točka.
        % \end{solution}

        \part[3]
        $c^7+343d^3c^4$

        \begin{solution}[3cm]
            $$m^7+27g^3m^4=m^4(m^3+27g^3)=m^4(m+3g)(m^2-3mg+9g^2)$$
            Za pravilno izpostavljanje skupnega faktorja $1$ točka, uporaba formule vsote kubov $1$ točka, pravilen končen rezultat $1$ točka.
        \end{solution}





    \end{parts}





      \newpage



        % 6. vprašanje
    \question
    Določite neznani števki, da bo:
    \begin{parts}
        \part[2]
        število $\overline{798312a456aa}$ deljivo s $4$;

        \begin{solution}[5cm]
            Za navedbo kriterija za deljivost z $9$ $1$ točka, uporaba kriterija in izračun $a\in\{0,9\}$ $1$ točka.
        \end{solution}


        \part[4]
        število $\overline{4x321xy}$ deljivo z $18$.

        \begin{solution}[6cm]
            Za navedbo kriterija deljivosti s $20$ $1$ točka, pravilno uporabljen kriterij za deljivost s $5$ $1$ točka, pravilno uporabljen kriterij za deljivost s $4$ $1$ točka, pravilen zaključek in zapis kombinacij $(x,y)\in\left\{(0,0),(2,0),(4,0),(6,0),(8,0)\right\}$ $1$ točka.
        \end{solution}


    \end{parts}

    \question[3]
    Dokažite: Razlika kuba poljubnega lihega števila in kvadrata njemu naslednjega sodega števila je liho število.

    \begin{solution}[7cm]
        Zapis $2\mid((2n-1)^2+(2n+1)^3)$ oziroma podobno $1$ točka; izračun kvadrata in kuba binomov in seštevanje $1$ točka; izpostavljanje $2$ $1$ točka, pravilen zaključek $1$ točka.
        
    \end{solution}



    % % \newpage
    % %dodatno vprašanje
    % \bonusquestion[1]
    % \textit{Dodatna naloga:} Dopolnite izraz do popolnega kvadrata in ga razcepite: $a^4-20a^2$.

    % \begin{solution}[5cm]
    %     $$a^4-20a^2+100=(a^2-10)^2$$
    %     Za pravilno dopolnitev $1$ točka, za pravilno zapisan razcep $1$ točka.
    % \end{solution}




    % % \newpage
        

    % % dodatno vprašanje
    % \bonusquestion[3]
    % \textit{Dodatna naloga:} Pokažite, da velja $77\mid 36 960$.

    % \begin{solution}[7cm]
    %     Za pravilno upoštevan kriterij za deljivost s $77$ $1$ točka, za pravilno uporabljen kriterij za deljivost s $7$ $1$ točka, za pravilno uporabljen kriterij za deljivost z $11$ $1$ točka.
    % \end{solution}




\end{questions}



\end{document}